\documentclass[11pt]{article}

% Packages
\usepackage{kotex}
\usepackage{graphicx}
\usepackage{indentfirst}
\usepackage{setspace}
\usepackage[linguistics]{forest}       % Syntax Tree 패키지 
\usepackage[a4paper, 
            top = 2.5cm, 
            bottom = 2.5cm, 
            left = 2cm, 
            right = 2cm]{geometry}     % 용지 여백 조절 패키지 
\usepackage{fancyhdr}                  % 머리글, 바닥글 조정 패키지
\pagestyle{fancy}
\fancyhf{} % 기존 머리글과 바닥글 초기화
\fancyhead[R]{몽골어 1: lab midterm quiz}   % 오른쪽: 문서 제목
\fancyfoot[L]{\nouppercase{\leftmark}} % 왼쪽: 현재 section 제목
\fancyfoot[R]{\thepage}                % 오른쪽: 페이지 번호
\renewcommand{\sectionmark}[1]{%
  \markboth{#1}{}%
}
\renewcommand{\headrulewidth}{0.4pt}   % 머리글 아래쪽 선
\renewcommand{\footrulewidth}{0.4pt}   % 바닥글 위쪽 선
\setlength{\headheight}{14pt}          % 머리글 높이 확보

% 제목 페이지용 바닥글
\fancypagestyle{firstpage}{%
  \fancyhf{}%
  \fancyfoot[L]{1. Tree}%
  \fancyfoot[R]{\thepage}%
  \renewcommand{\headrulewidth}{0pt}%
  \renewcommand{\footrulewidth}{0.4pt}%
}

\setcounter{secnumdepth}{0}            % section 넘버링 제거
\forestset{
  terminal/.style={
    no edge,
    l=3ex,
    for parent={l sep=0pt}
  }
}

% Short form for a lexical category and its terminal.
\forestset{
  lex/.style={append={[#1, terminal]}},
  simple terminals/.style={
    delay={for tree={split option={content}{ }{content,lex}}}
  }
}

% Reusable GPS tree block; default width is 0.48\textwidth.
\NewDocumentEnvironment{gpstree}{O{0.48\textwidth} m +b}
{%
  \begin{minipage}[t]{#1}
  \vspace{0pt}%
  \textbf{#2}\par
  \medskip
  \centering
  \makebox[\linewidth][c]{%
    \begin{forest} simple terminals
      #3
    \end{forest}%
  }
  \end{minipage}%
}
{}

% Metadata
\title{2026 Fall 몽골어 1(001): \\ \medskip \large{laboratory midterm quiz}}
\author{Болорцэцэг  /  1-102}
\date{October 13, 2026}

\begin{document}
\maketitle

\section{1. 제시된 자료를 보고, 알맞은 답을 적으시오.}
\bigskip
\noindent
\begin{tabular}{@{}p{0.47\textwidth}p{0.47\textwidth}@{}}
  1-1. долоo \textbf{(답: 7)} &
  1-4. 56 \textbf{(답: Тавин зургаа)} \\
  [0.8em]
  1-2. арван гурав \textbf{(답: 13)} &
  1-5. 74 \textbf{(답: Далан дөрөв)} \\
  [0.8em]
  1-3. Гучин долоо \textbf{(답: 37)} &
  1-6. 99 \textbf{(답: Ерэн ес)}
\end{tabular}
\bigskip

\section{2. 제시된 한국어 뜻을 보고 그에 해당하는 몽골어 단어를 적거나, 몽골어 단어를 보고 그것의 한국어 뜻을 적으시오.}
\bigskip
\noindent
\begin{tabular}{@{}p{0.47\textwidth}p{0.47\textwidth}@{}}
  2-1. 여권 \textbf{(답: паспорт)} &
  2-4. гэр \textbf{(답: 게르, 집)} \\
  [0.8em]
  2-2. 학생 \textbf{(답: оюутан)} &
  2-5. курс \textbf{(답: (대학교의) 학년)} \\
  [0.8em]
  2-3. 가족 \textbf{(답: гэр бүл)} &
  2-6. найз \textbf{(답: 친구)}
\end{tabular}
\bigskip  

\section[3. 십자말풀이]{3. 십자말풀이\footnotemark}
\footnotetext{십자말풀이 표 모양을 정확히 복원하기 어려워 십자말풀이에 제시된 설명만 기재하였음.}
\bigskip
\noindent
\begin{tabular}{@{}p{0.47\textwidth}p{0.47\textwidth}@{}}
  3-1. Би \_\_\_(20살) настай. \textbf{(답: хорь)} &
  3-4. 3 + 4 = ? \textbf{(답: долоo)} \\
  [0.8em]
  3-2. 5 + 6 = ? \textbf{(답: арван нэг)} &
  3-5. Миний ээж \_\_\_(50살) настай. \textbf{(답: тавин)} \\
  [0.8em]
  3-3. Би \_\_\_ нэгэн(21살) настай. \textbf{(답: xорин)} &
  3-6. 2 + 2 = ? \textbf{(답: дөрөв)}
\end{tabular}
\bigskip

\section{4. 자기소개: 몽골어로 자신을 소개할 수 있는 문장을 적으시오. (2문장 이상)}
\medskip
\noindent \textbf{예시 답안:} Миний монгол нэр Баатар. Би Сөүлийн Их Сургуулийн оюутан. Би хоёрдугаар(2-p) курсийн оюутан. 
(내 몽골 이름은 Баатар입니다. 나는 서울대학교 학생입니다. 나는 대학교 2학년입니다.)

\end{document}